\section{Implementation}

\subsection{Implementation Environment}
\label{sec:impl-env}

The implementation is based on gem5 version 23.0.0.1, and was developed on a separate branch. The simulator is built for the RISC-V ISA. All experiments use
syscall-emulation mode with a single \texttt{RiscvMinorCPU} that has
one hardware thread. The workloads are statically linked
RISC-V programs built for the \texttt{rv64g} ISA.

The baseline configuration models a 1\,GHz \texttt{RiscvMinorCPU}
with timing memory accesses and 64-byte cache lines. The L1
instruction cache is 32\,KiB and 4-way. The L1 data cache is 64\,KiB
and 4-way, with 16 MSHRs and a stride prefetcher. The L2 cache is
256\,KiB and 8-way, with 20 MSHRs and an AMPM prefetcher, and main
memory is DDR4-2400. All pipeline widths are four. At most one memory
instruction is issued and one is committed per cycle, and early memory
issue is enabled. The CPU uses the \texttt{SingleThreaded} thread
policy with up to 16 outstanding memory accesses.

\subsection{LVP Integration}
\label{sec:lvp-integration}

The LVP predictor is integrated into the Execute stage of gem5
MinorCPU. Execute
must link each prediction to the correct dynamic producer and decide
when a dependent load may use it. It must also substitute the
predicted operand while generating the memory request, and validate
the prediction when the producer's actual response arrives. If an
incorrect prediction has influenced speculative execution, the
younger instruction stream must be discarded and fetched again.

The implementation described here covers standalone LVP-based
load-to-load address speculation.

\subsubsection{Configuration and Prediction Eligibility}

Three configuration options decide whether LVP-based address
speculation can operate. \texttt{enableLVP} enables the predictor,
and \texttt{lvpSpecAddr} lets dependent loads use confident
predictions as their base addresses. Address speculation also needs
\texttt{executeAllowEarlyMemoryIssue} to be enabled. Invalid
combinations are rejected when the Execute stage is constructed.

An optional PC range can restrict address speculation to part of the
program. If no range is set, speculation is allowed everywhere.

By default, the predictor uses a 1,024-entry, untagged table
indexed by the load instruction's PC:

\[
    \operatorname{index}(PC)
    = (PC \gg 1) \mathbin{\&} (N-1),
    \qquad N=1024.
\]

Each entry stores the most recently observed value and a 2-bit
saturating confidence counter. The counter ranges from zero to
three. The default prediction threshold is three, and it can be
configured from one to three. A lookup returns a prediction only
when the entry's confidence reaches the configured threshold. When
the observed value matches the stored one, the counter is
incremented, up to its maximum. When it does not match, the counter
is decremented, down to zero. In both cases, the stored value is
replaced by the newly observed value. Entries start with value zero
and confidence zero. The table has no separate valid bit and no PC
tag.

The implementation also restricts which instructions can act as
producers and consumers. An eligible producer is an ordinary integer
load with one integer destination register, excluding the zero
register. Atomics, store-conditionals, micro-operations, serializing
instructions, and prefetches are excluded here. An eligible consumer must
meet the same ordinary-load restrictions and have exactly one
integer source register, which is its base address.

As a result, the speculative operand path is limited to dependent
loads. General ALU instructions do not directly consume LVP
predictions.


\subsubsection{Issue-Time Prediction and Dependency Tracking}

Prediction lookup and dependency tracking run when an instruction
issues into a functional unit, in program order.

The predictor table and \texttt{lvpSpecMap} serve different purposes.
The predictor table stores learned values indexed by static
instruction PCs. \texttt{lvpSpecMap}, on the other hand, tracks the
predicted writer that is currently available for each architectural
register. It is kept separately for each thread. It maps a destination
register to an entry containing the producer's dynamic
sequence number, PC, and predicted value:

\[
\texttt{lvpSpecMap[tid][reg]}
=
\{\texttt{producerSeq},\texttt{producerPC},\texttt{value}\}.
\]

Issue-time processing has three steps, done in this order. The
predictor lookup in the third step happens whenever LVP is enabled.
The consumer lookup, destination invalidation, and insertion into
\texttt{lvpSpecMap} happen only when address speculation
(\texttt{lvpSpecAddr}) is also enabled. Without it, predictions are
made and later validated, but no consumer can use them.

\paragraph{1. Consumer dependency lookup:}
If the issued instruction is an eligible consumer, Execute looks at
its base register and searches \texttt{lvpSpecMap} for a matching
entry. Finding an entry is not enough. The recorded producer sequence
number must equal the consumer's \texttt{instToWaitFor} value, which
identifies the dynamic instruction the consumer is waiting on. Both
the consumer PC and the recorded producer PC must also fall inside
the configured speculation window.

When these checks pass, the consumer saves the producer's identity
and the predicted value with itself. This snapshot keeps the
producer-consumer relationship intact even if the map is updated
later.

\paragraph{2. Destination invalidation:}
Execute then removes the existing map entries for every destination
register written by the current instruction. This stops an older writer's prediction from staying available after a newer instruction
has become the latest writer of that register. The invalidation is not
limited to eligible loads. Any instruction that writes a register
invalidates the corresponding entry.

\paragraph{3. Producer prediction lookup:}
Finally, if the current instruction is an eligible producer,
Execute looks up its PC in the LVP table. If the lookup returns a
confident prediction, the dynamic instruction is marked as predicted,
and its predicted value and producer state are recorded. When address
speculation is enabled, the prediction is also inserted into
\texttt{lvpSpecMap} under its destination register. The
entry identifies this particular dynamic instance by its execution
sequence number.

The sequence number is needed because several dynamic instructions
can write the same architectural register. A consumer must use the
prediction of the exact producer named in its dependency information,
not just any prediction attached to that register.

\subsubsection{Admission for Early Memory Issue}

Saving a prediction in the consumer does not start speculative
execution by itself. The unresolved dependency is the
predicted load at the head of the in-flight queue. This path is tried
only when the ordinary early-issue path cannot admit the consumer.

The consumer is admitted through LVP only when the oldest memory reference is waiting to enter the LSQ. It
must still hold a saved prediction whose producer matches its
\texttt{instToWaitFor}, and it must not have been accepted yet. Its
producer must be the in-flight head, still be marked as predicted,
already be in the LSQ, and not yet have its memory response. Both
instructions must belong to the current instruction stream.

The producer must be the head, so every older instruction has already
left the in-flight queue and the consumer never skips past a chain of
unresolved dependencies. The producer must also be in the LSQ without
a response, so speculation happens only while the real load result is
still pending.

\subsubsection{Operand Substitution and LSQ Request Generation}

Once the consumer is admitted, Execute turns on operand substitution
for the duration of the call that issues the load. While it is on, a
read of the consumer's base register returns the saved predicted value
instead of the register file value. The substitution applies only to
that one operand of that one instruction, and the register file
itself is never changed.

The load then forms its effective address from the predicted base and
sends a real memory request to the LSQ through the normal path. Only
when this request has actually entered the LSQ is the consumer marked
as accepted and the producer marked \texttt{lvpUsedBySpec}. Saving a
prediction or passing the admission checks is not enough. Only an
accepted request counts as speculative use, and only then can a later
misprediction of the producer require recovery.

\subsubsection{Producer Validation and Misprediction Recovery}

The prediction is checked when the producer's memory response
arrives. Execute handles completed instructions in program order, so
the producer is always checked before any younger consumer can
commit. Once the real value has been written to the destination
register, the LVP table is trained with it. This happens for every
eligible load, whether or not it was predicted.

If the producer was predicted, the real value is compared with the
prediction. Recovery is needed only when the prediction was wrong and
the producer is marked \texttt{lvpUsedBySpec}. A wrong prediction
that no accepted consumer used causes no redirect.

For recovery, Execute keeps the producer and redirects fetch to the
instruction right after it. This starts a new instruction stream,
which clears \texttt{lvpSpecMap}. Every younger instruction from the
old stream is then discarded by MinorCPU's normal stream check, not
only the ones that used the wrong value, and is fetched again. The
same check stops a late memory response from the old stream from
completing. As an extra safeguard, an accepted consumer checks, when
its own response arrives, that its base register now holds the value
it used. By then the producer has committed, so a mismatch can only
mean a recovery bug, and the simulation stops.

The redirect also needs a small change in Fetch2. The younger branches may already have updated the
branch predictor's speculative history. Fetch2 recognizes this case
and rolls the branch predictor back to the redirect point. It does
not treat the load as a mispredicted branch.

Each instruction also tracks its producer and consumer roles
separately, because a dependent load can use an older prediction
while its own value is also predicted. When an instruction leaves the
pipeline, the outcome of each role (committed, squashed, or faulted)
is recorded. A producer's \texttt{lvpSpecMap} entry is removed only
if it still belongs to that same producer, so cleanup never deletes a
newer writer's prediction.

\subsection{ASTRA Address Prediction}
\label{sec:impl-astra}

ASTRA is a PC-indexed predictor made of a Last Address Table (LAT)
and a Stride$\rightarrow$Stride table (SS), with lookup and training
done in the Execute stage. The paper's plain ASTRA stores the last
address of a load. This implementation instead stores the load's last
base-register value and predicts the next base value. The instruction's displacement is added only
later, when V-ASTRA forms an effective address
(Section~\ref{sec:impl-vcache}). Addresses keep their full 64-bit
width.

\subsubsection{Predictor Organization}

Each LAT entry holds a 64-bit base and an index into the SS table,
and each SS entry holds a signed 16-bit stride. Both tables have \(N\)
entries, 1,024 by default. A load's LAT row \(r\) is chosen from its
PC, and a stride \(S\) selects SS entry \(S \bmod N\), that is, its
lower bits:
\[
    r = (PC \gg 1) \mathbin{\&} (N-1).
\]
The tables are untagged and have no valid bits or confidence
counters. They start filled with zeros, so every eligible load gets a
prediction.

\subsubsection{Prediction and Base Capture}

ASTRA makes a prediction when an eligible load issues. Eligible loads
are ordinary loads with an integer base register, with the same
exclusions as LVP. ASTRA reads the load's LAT entry to get the
last stored base \(L\) and its SS index \(k\). It then adds the
stride stored there, sign-extended from 16 bits, to get the predicted
base:
\[
    B = L + SS[k].
\]

Just before the load sends its memory request, the base it actually
uses is captured for training and for the V-ASTRA check. Normally,
this is the value in the base register. If the base came from a
validated V-ASTRA value through forwarding
(Section~\ref{sec:impl-sb}), the forwarded value is captured instead.
If it came from an LVP prediction, nothing is captured by default,
because that value is not yet known to be correct. The experimental
modes of Section~\ref{sec:impl-provisional} change this.

On its own, ASTRA only makes and records predictions. It does not
form an effective address or skip the memory access. That is done by
V-ASTRA (Section~\ref{sec:impl-vcache}).

\subsubsection{Training}

ASTRA is trained when a load retires, so prediction happens at issue
and the update happens at retirement. For a load whose captured base
is \(A\), the new stride is
\[
    S = A - L.
\]
The stride is stored in 16 bits, so larger differences wrap around.
The SS entry used for the last prediction is updated first, and then
the LAT records the new base and the index of the new stride:
\[
    SS[k] \leftarrow S, \qquad
    LAT[r] \leftarrow \bigl(A,\ S \bmod N\bigr).
\]
Only loads that commit with a captured base train ASTRA. Squashed and
faulting loads do not. Training waits for retirement, so several
instances of the same load may issue before the first one updates the
tables, and these later instances use the older state. No speculative
updates are made to hide this delay.

\subsection{Value Cache and Load Bypass}
\label{sec:impl-vcache}

V-ASTRA adds a separate value cache and a guarded load bypass on top
of ASTRA. It works in two stages. Execute first decides whether a load
is a candidate, and the LSQ later decides whether the bypass actually
happens. Being a candidate does not guarantee that the load will skip
memory.

\subsubsection{Value Cache Organization and Updates}

The value cache is a 4\,KiB direct-mapped structure with 256 entries
of 16 bytes each. Its index and tag come from the virtual address:
\[
    \operatorname{index}(A) = (A \gg 4) \mathbin{\&} 255,
    \qquad
    \operatorname{tag}(A) = A \gg 12.
\]
Each entry has an occupancy bit, a tag, 16 data bytes, and one valid
bit per byte. A lookup hits only if the entry is occupied, the tag
matches, and every byte the load needs is valid.

The cache is updated when memory operations complete, as long as they
did not fault and were not discarded. A load fills or replaces the
entry at its index, but only the bytes it actually read are copied in
and marked valid. Neighboring bytes are not filled. A store updates an
entry only if the tag already matches, and never allocates a new one.
Atomics, load-reserved and store-conditional operations, uncacheable
accesses, and other special accesses invalidate the bytes they touch.
Loads that cross a 16-byte boundary are ignored, while stores that
cross one invalidate the affected bytes in both entries. A bypassed
load does not update the cache, because its data did not come from
memory. The whole cache is also flushed after every system call
(\texttt{ecall}), since the operating system may have written memory.

\subsubsection{V-ASTRA Candidate Selection}

Candidate selection happens in Execute, right after the actual base
has been captured and before the normal memory access starts. First,
ASTRA's predicted base is compared with the actual base. If they
differ, or no base was captured, the load is not a candidate. If they
match, the effective address is formed from the predicted base \(B\)
and the instruction's signed 12-bit displacement \(D\):
\[
    EA = B + D.
\]
The load is then checked against older stores. Issued stores,
atomics, and store-conditionals are kept in a per-thread list with
their addresses. The load's exact size is not known yet, so the scan
conservatively checks an 8-byte range. An overlapping older store
rejects the candidate, and so does an older store whose address is
still unknown. The second case should not happen, because memory
instructions enter the LSQ in program order, but the check is kept as
a safeguard.

If the load passes, it becomes a candidate. Its effective address and
a snapshot of the matching value-cache entry (tag match, valid bytes,
and data) are saved for the later bypass decision. A candidate does
not need a value-cache hit at this point. That is checked later in
the LSQ.

\subsubsection{LSQ Bypass and Safety Checks}

The actual bypass is decided later, when the load's request is about
to leave the LSQ for memory. By then the address has been translated,
so the TLB has already been accessed. The LSQ first applies its usual
rules. The request must be in the current stream, have no translation
fault, and not need load-reserved/store-conditional handling. No older
store may still be waiting ahead of it. If the store buffer fully
covers the load, normal store-to-load forwarding is used instead. If
it only partly covers the load, the load waits.

The bypass requires the load to be a candidate that has not already fallen back to memory, and it must
be a plain single read that fits inside one 16-byte value-cache entry.
Its real address must equal the address saved at candidate selection.
Older stores are scanned again, this time with the exact access size.
The saved snapshot must have a matching tag and all requested bytes
valid. A fresh value-cache lookup must still return the same
bytes, in case the cache changed after the candidate was chosen.

If any check fails, the load falls back to a normal memory access. If all checks pass, the value-cache bytes
are copied into the load's response and the request is marked
complete. No memory request is sent. The address translation has
already happened, so the bypass saves the memory access, not the
translation. The load then completes and commits in order like any
other load, and its value is written to the destination register at
commit. Its data did not come from memory, so it does not update the
value cache.

\subsubsection{Strict Verification Mode}

A strict verification mode was added to check the value cache's data
without relying on it. In this mode, a candidate whose snapshot has a
matching tag and valid bytes still goes to memory normally. When the
real data returns, it is compared with the saved snapshot. A match is
counted, and a mismatch stops the simulation with the details of the
load. This mode cannot be combined with the actual bypass, because a
bypassed load never receives real memory data to compare against.

\subsection{Scoreboard Release and Dependent-Load Forwarding}
\label{sec:impl-sb}

MinorCPU's scoreboard tracks when each register's value is expected
to be ready and which in-flight instruction will write it. It does
not hold the values themselves. Instructions read their operands from
the register file when they execute, except when the LVP substitution
or the forwarding described below is active.

The dependent wait depends on \texttt{extraAssumedLat}. With
the default value of 2, a load's destination is not marked
unpredictable, so dependents may issue before the load commits,
although they can still be held back until it does. To model a strict
stall-on-use policy, a separate configuration sets
\texttt{extraAssumedLat} to zero. A load's destination is then marked
unpredictable, and dependents cannot issue until that mark is cleared.

\subsubsection{Design~A: Scoreboard Release}

Design~A lets the dependents of a bypassed load become eligible to
issue before the load commits. Under stall-on-use timing, an issued
load's destination registers are marked unpredictable. Once the
V-ASTRA bypass has succeeded and the load's response has been
created, its result is known. At that point, Design~A clears the
mark. Nothing else in the scoreboard or register file changes. A flag
records that the mark was cleared, so
it is not cleared a second time when the load later commits or is
discarded. If the bypass fails, nothing is released and the load
continues normally.

This only changes when a dependent may issue. It does not give the
dependent the producer's value, which is still read from the register
file once the producer commits. A dependent can enter its functional
unit earlier, but a dependent load still cannot send its memory
request early through Design~A alone. Under the default memory-latency
setting, loads are never marked unpredictable, so the release has no
effect.

\subsubsection{Design~B: Dependent-Load Forwarding}

Design~B goes further and forwards the validated value to a dependent
load. For each V-ASTRA candidate, a forwarding record is created with
the producer's destination register, stream, value, and state, unless
the candidate's base is a provisional LVP value
(Section~\ref{sec:impl-provisional}). The record starts as
\texttt{Pending}, because being a candidate does not mean the bypass
will succeed. It becomes \texttt{Validated} only when the bypass
succeeds, and \texttt{Failed} if the load falls back to memory or
cannot be forwarded.

When the record is validated, the value is rebuilt from the checked
value-cache bytes, with the right sign or zero extension for the
integer loads \texttt{lb}, \texttt{lbu}, \texttt{lh}, \texttt{lhu},
\texttt{lw}, \texttt{lwu}, and \texttt{ld}. Floating-point loads
cannot be producers.

A dependent load can use the forwarded value only if it is an ordinary
load with a single integer source which is at the front of the memory
functional unit and is waiting on exactly this producer
(\texttt{instToWaitFor}) in the same stream, with matching registers.
The record must be \texttt{Validated}. With a pending or failed
record, the consumer simply waits.

When admitted, the consumer reads the forwarded value instead of its
base register, using the same kind of temporary substitution as LVP.
The register file itself is not changed. It forms its address and sends
its memory request before the producer commits. ASTRA also captures
this forwarded value as the consumer's base, so a forwarded load can
itself become a candidate, bypass, and forward to the next load. This
allows chains of forwarded loads. When the consumer's own response
arrives, the producer has committed, so the register must now hold
the forwarded value. A mismatch stops the simulation. Records are
removed when instructions leave the pipeline, and checks at drain time
confirm that none are left over.

Forwarding needs V-ASTRA bypass, scoreboard release, and early memory
issue to be enabled. It is off by default and supports a single
hardware thread, like all value-cache features. I used it with the
\texttt{SingleThreaded} commit policy. Without META, it cannot be
combined with LVP address speculation. Neither design forwards values
to arithmetic instructions, and neither is out-of-order execution.

\subsection{META Integration}
\label{sec:impl-meta}

META sits in MinorCPU's early memory-issue path. When a dependent load
at the front of the memory functional unit wants to issue before its
producer commits, META decides whether it may use the V-ASTRA value,
the LVP prediction, or neither. The decision is repeated on every
commit pass until the load is admitted or can proceed normally.

This differs from the paper's META, which picks one prediction per
load at issue time. Here, META chooses the source of a dependent
load's base operand. A choice only lets the load issue early in that
cycle, and the source is recorded only if the load's memory request
actually enters the LSQ. A load that is never admitted early counts
as having used no prediction.

\subsubsection{Selection Policy}

META first finds the producer the dependent load is waiting on and
looks up that producer's V-ASTRA forwarding record. It then applies
this priority:

\begin{enumerate}
    \item \textbf{Validated V-ASTRA value:} if the record is
    validated and the forwarding conditions hold, V-ASTRA is used.
    \item \textbf{Pending V-ASTRA outcome:} if the record is still
    pending, the load waits, even if LVP could serve it. META does not
    fall back to LVP while the V-ASTRA outcome is undecided.
    \item \textbf{LVP prediction:} if the record failed or does not
    exist, LVP is used if its normal admission conditions hold.
    \item \textbf{No prediction:} otherwise, the load waits for its
    producer as usual.
\end{enumerate}

META reuses the existing V-ASTRA and LVP admission checks and does not
change either one. It only decides which of them, if any, may serve
the load. If both could serve it, V-ASTRA is chosen, and a check
makes sure that a load is never accepted through both paths.

\subsubsection{Source Attribution and Recovery}

When the load enters the LSQ, the chosen source is recorded with it.
The two sources need different recovery, so this record is important.
An LVP-served load used a speculative value, so if the producer's
prediction turns out wrong, the normal LVP recovery applies. A
V-ASTRA-served load used an already validated value, so it never
causes LVP recovery, even if the producer also had a wrong LVP
prediction. That prediction is still checked and used to train LVP as
usual.

When a load leaves the
pipeline, META records its source and whether it committed, was
squashed, or faulted. It also checks that no load was accepted
through both sources.

\subsubsection{Configuration}

META is off by default. Turning it on requires LVP with address
speculation and V-ASTRA forwarding, which in turn need ASTRA, the
value cache, V-ASTRA bypass, scoreboard release, and early memory
issue.

\subsection{Experimental ASTRA Interaction for LVP-Served Loads}
\label{sec:impl-provisional}

By default, a load whose base comes from an LVP prediction gives ASTRA
nothing. Its base is not captured, so it does not train ASTRA and
cannot become a V-ASTRA candidate. This limits how the two predictors
work together, especially on workloads where LVP serves many loads.

To study this, I added the \texttt{metaAstraLvpLoads} option with
three modes. The default, \texttt{off}, keeps the behavior just
described. \texttt{train} lets the LVP-supplied base train ASTRA, but
only after it has been validated. \texttt{candidate} goes further and
also lets that base be used for V-ASTRA candidate selection right
away, before it is validated. None of the modes changes how META
chooses between LVP and V-ASTRA.

\subsubsection{Provisional Base Capture and Validation}
\label{sec:impl-provisional-capture}

In the \texttt{train} and \texttt{candidate} modes, the LVP-supplied
base of a load is captured as a provisional base. It is recorded, but
marked as not yet validated. Two separate flags are kept. One says
whether the base is provisional. The other says whether the base may
be used for the V-ASTRA decision. In \texttt{candidate} mode, both are
set at once. A base can then be used for candidate selection while
it is still unvalidated.

The base is validated when the load's own memory response arrives.
At that point, the base register is compared with the LVP value the
load used. If they match, the provisional mark is cleared, and the
base is replaced with the register value. The load then trains ASTRA
at retirement in the normal way. This check happens before any fault
on the response is handled, so a load with a correct base but a
translation fault is counted as validated and then retires as faulted,
without training ASTRA.

If the load is discarded, or is re-executed after an early fault
before being validated, its provisional base is dropped and never
used for training. A load is never allowed to commit while its base
is still provisional.

\subsubsection{Mode \texttt{train}: Deferred ASTRA Training}

In \texttt{train} mode, the LVP-supplied base is captured as
provisional but is not usable for the V-ASTRA decision. The load
therefore cannot become a candidate or bypass memory. When its
response arrives and the base is validated, the base is replaced with
the register value, and the load trains ASTRA when it retires. If the
LVP prediction was wrong, normal LVP recovery discards the load, and
its provisional base is dropped. A wrong LVP value can therefore never
become an ASTRA training value. This mode gives ASTRA more
information without letting an unvalidated base affect any V-ASTRA
decision.

\subsubsection{Mode \texttt{candidate}: Provisional V-ASTRA Candidacy}

The \texttt{candidate} mode is more aggressive. The LVP-supplied base
is captured as provisional and is also usable for the V-ASTRA
decision right away. V-ASTRA compares ASTRA's predicted base \(B\)
with this LVP value. If they match, the effective address is formed
as usual:
\[
    EA = B + D.
\]
A match only shows that ASTRA and LVP agree. It does not show that
the LVP value is correct.

The candidate must still pass all the normal V-ASTRA safety checks i.e. the older-store scans, the snapshot check, and the fresh value-cache
lookup in the LSQ. The address checked there is the one computed from
the LVP value, so it confirms agreement with ASTRA. A bypass can therefore happen before the LVP prediction
is validated. This is safe because the load is already an accepted
LVP consumer. Its result is only a pending response. If the LVP
prediction turns out wrong, normal LVP recovery discards the load
before it commits. If the prediction is validated, the load trains
ASTRA at retirement. Like other bypassed loads, it does not update the
value cache.

A load with a provisional base may not create a forwarding record.
Otherwise, a younger load could receive a value derived from an
unvalidated base. Such dependents simply wait for the producer as
usual. The forced-misprediction tests check that every such bypassed load follows the order below, where each
event happens in the same cycle as the previous one or later:
\[
    \text{bypass}
    \;\rightarrow\;
    \text{validation}
    \;\rightarrow\;
    \text{redirect}
    \;\rightarrow\;
    \text{discard}.
\]

\subsubsection{Fault Handling and Provisional-State Retry}

One fault case exposed a bug in the first version of this mechanism.
An early-issued load can use an LVP value to form its address and then
hit an alignment fault before its request enters the LSQ. On RISC-V,
this happens only when alignment checking is enabled. MinorCPU then
stops the early attempt and runs the load again once it reaches the
head of the in-flight queue. A translation fault is different,
because it is reported with the response, after the provisional base
has already been validated (Section~\ref{sec:impl-provisional-capture}).

When the load runs again, the LVP substitution from the first attempt
is no longer active, so the base is taken from the register file.
Before the fix, the provisional mark from the first attempt could
remain. The load then reached commit with a stale provisional base,
which triggered the rule that a committed load must not have one. The
fix drops any leftover provisional base at the start of every new
base-capture attempt. Each attempt now starts clean.

In \texttt{candidate} mode, the first attempt's V-ASTRA decision is
also kept across the fault. On re-execution, the real address may
differ from the one saved earlier, and the later bypass check then
rejects the request because the two addresses do not match. If the LVP
value had been correct, the re-executed load would use the same
misaligned address and fault again, so it would never reach the
bypass point.

\subsubsection{Configuration}

The \texttt{metaAstraLvpLoads} option takes the values \texttt{off},
\texttt{train}, and \texttt{candidate}, with \texttt{off} as the
default. The two experimental modes require META to be enabled.
Counters record how many provisional bases are created, validated,
and dropped. Validation and performance results for the three modes
are given in Sections~\ref{subsec:meta-integration-recovery-validation}
and~\ref{subsec:provisional-base-performance}.
